
\section{Introduction}\label{sec:intro}

\Glspl*{lvm} are ubiquitous in many areas of statistical science, e.g., complex probabilistic models for text, audio, video and images \citep{blei2003latent, smaragdis2006probabilistic, hoff2002latent}, or inverse problems \citep{glyn2024statistical}. These models capture the underlying structure of the data in terms of low-dimensional variables -- a structure that often exists in the real world \citep{whiteley2022statistical}. However, \glspl*{lvm} often require non-trivial procedures for their estimation as their likelihood is often intractable.

Consider a generic latent variable model as $p_\theta (x,y)$, parameterised by $\theta\in\mathbb{R}^{d_\theta}$, for \textit{fixed}, observed data $y \in \bR^{d_y}$, and latent variables $x \in \bR^{d_x}$. Thus, formally, we see the statistical model as a real-valued mapping $p_\theta(x,y) :\mathbb{R}^{d_x}\times \mathbb{R}^{d_\theta}\to \bR$. The task we are interested in is to estimate the parameter $\theta$ that explains the fixed dataset $y$. Often, this is achieved via the \gls*{mle}. In our setting, due to the presence of latent variables, we aim at finding the \gls*{mmle}, which is termed the \gls*{mmle} problem \citep{dempster1977maximum}. More precisely, our problem takes the form
\begin{align}\label{eq:max}
\bar{\theta}_\star \in \argmax_{\theta\in\mathbb{R}^{d_\theta}} \log p_\theta(y),
\end{align}
where $p_\theta(y) := \int p_\theta(x,y) \mathrm{d}x$ is the \textit{marginal likelihood} (also called the model evidence in Bayesian statistics \citep{bernardo2009bayesian}). It is apparent from \eqref{eq:max} that the problem cannot be solved via optimisation techniques alone for most statistical models.

Classically, this problem is solved with the iterative \gls*{em} algorithm \citep{dempster1977maximum}, which increases the marginal likelihood at every iteration and, under regularity conditions, converges to a stationary point of $\theta \mapsto \log p_\theta(y)$ \citep{wu1983convergence}. At a given parameter estimate $\theta_n$, the \gls*{em} algorithm implements a map $\theta_n \mapsto \argmax_{\theta \in \mathbb{R}^{d_\theta}} \mathbb{E}_{p_{\theta_n} (x|y)}[\log p_{\theta} (x, y)]$, which is guaranteed not to decrease the marginal likelihood at any iteration. This procedure requires an ``E-step'', namely the computation (or estimate) of the expectation w.r.t. the posterior distribution of the latent variables $p_{\theta_n}(x|y)$, and an ``M-step'', which is the maximisation of the expected log-likelihood w.r.t. $\theta$. These steps are intractable in general and can be approximated with a variety of methods, which have been extensively explored. For example, Monte Carlo \gls*{em} \citep{wei1990monte} and stochastic \gls*{em} \citep{celeux1985sem} have been widely studied  \citep{celeux1992stochastic, chan1995monte, sherman1999conditions, booth1999maximizing, cappe1999simulation, diebolt1995stochastic}.

In the most general case, the \gls*{em} algorithm is implemented using \gls*{mcmc} techniques for the E-step \citep{atchade2017}, and numerical optimisation techniques for the M-step \citep{meng1993maximum,liu1994ecme,lange1995gradient}. With the popularity of \gls*{ula} in Bayesian statistics and machine learning \citep{durmus2017nonasymptotic,durmus2019high,dalalyan2017theoretical}, novel variants of \gls*{em} algorithms use unadjusted chains. Most notably, \citep{de2021efficient} studied an algorithm termed \gls*{soul}, which performs unadjusted Langevin steps for the E-step and stochastic gradient ascent for the M-step, building on the ideas of \citep{atchade2017}. In the \gls*{mcmc} based \gls*{em} setting, the bias incurred \emph{by using finitely many} \gls*{mcmc} steps complicates the theoretical analysis and requires stringent conditions for step-sizes.

An alternative approach was developed in \citep{kuntz2023particle} where the authors proposed an interacting particle system, consisting of $N$ particles to replace sequential \gls*{mcmc} steps. This framework, termed \gls*{pgd}, is based on the idea of joint optimisation and sampling to avoid a double-loop structure as in \gls*{mcmc}-based approaches, which results in an efficient algorithmic framework as well as one that is amenable to theoretical analysis \citep{caprio2024error}. Inspired by the approach in \citep{kuntz2023particle}, a closely related interacting particle system was proposed in \citep{akyildiz2025interacting}, where a scaled noise is injected in the $\theta$-dimension. The authors termed this method \gls*{ipla} and proved error bounds for the algorithm. This seemingly small modification is significant, making the algorithm an instance of a Langevin diffusion (an observation we build on in this paper), streamlining the theoretical analysis and opening the door to a wide variety of extensions with theoretical guarantees. The \gls*{ipla} methodology has been already utilised in other contexts, see, \citep{johnston2024taming} for superlinear extensions and \citep{encinar2025proximal} for a set of proximal methods based on \gls*{ipla}, and \citep{akyildiz2024multiscale} for relations between \gls*{ipla} and multiscale methods. \\

\noindent\textbf{Contributions.} In this paper, we develop optimisation methods for \gls*{mmle} based on underdamped Langevin samplers. Our contributions can be summarised as follows: \\

\noindent\textbf{(C1)} We propose the \gls*{kipld}, a diffusion process to optimise the marginal likelihood in latent variable models. We prove that the stationary measure of this diffusion concentrates around the \gls*{mmle} (Propositions~\ref{prop:stationary} and \ref{prop:mmleerror}) and the \gls*{kipld} converges to this measure exponentially fast (Proposition~\ref{prop:kipld_conv}). \\
\glsunset{kiplmc1}
\glsunset{kiplmc2}

\noindent\textbf{(C2)} We then develop our first \gls*{kiplmc} method, termed \gls*{kiplmc1} -- an exponential integrator discretisation of \gls*{kipld}. We prove discretisation error bounds for our method showing convergence rates in both time and step-size (Theorem~\ref{thm:KPLMC1}), closing the problem of nonasymptotic analysis of this method in the strongly convex case. In particular, we show that \gls*{kiplmc1} attains accelerated rates of convergence compared to its overdamped counterparts \gls*{ipla} \citep{akyildiz2025interacting} and \gls*{pgd} \citep{caprio2024error}. While these overdamped diffusion-based \gls*{mmle} methods achieve an $\varepsilon$ error in Wasserstein-2 distance in $\widetilde{\mathcal{O}}(d_x \varepsilon^{-2})$ steps, we prove that \gls*{kiplmc1} attains $\varepsilon$ error in $\widetilde{\mathcal{O}}(\sqrt{d_x} \varepsilon^{-1})$ steps.\\

\noindent\textbf{(C3)} We then introduce a splitting-based explicit discretisation scheme within this setting, which leads to a novel \gls*{mmle} algorithm, termed \gls*{kiplmc2}. We study the nonasymptotic behaviour of this method (Theorem~\ref{thm:KIPLMC2}), showing that the algorithm obtains an error $\varepsilon$ in Wasserstein-2 distance in $\widetilde{\mathcal{O}}(d_\theta\sqrt{d_x}\,\varepsilon^{-3})$ steps.\\

\noindent\textbf{(C4)} Finally, we provide solid empirical evidence of the performance of our methods in a variety of settings, including synthetic and realistic models. In the convex setting, we validate the $\mathcal{O}(N^{-1/2})$ concentration onto the maximiser $\bar{\theta}_\star$  and demonstrate the increased stability of the momentum-based methods. In the non-convex setting, we show the acceleration of the proposed methods, as well as the importance of the choice of numerical integrator, validating our methodological contribution. \\

\noindent\textbf{Relations to existing work.} The \gls*{kipld} is similar to the continuous time system on which the \gls*{mpgd} algorithm by \citep{lim2023momentum} is based (see Eq.~(21) in \citep{lim2023momentum} where we take $\eta_\theta = \eta_x = 1$ and $\gamma_x = \gamma_\theta = \gamma$). In line with the work done by \citep{akyildiz2025interacting}, we inject the $\theta$-dynamics with an appropriately scaled Brownian motion, i.e., $\sqrt{2\gamma/N} \md \B^0_t$. The advantage of studying and implementing this version is due to the fact that the new system can be shown to be an example of a standard underdamped Langevin diffusion, which enables us to show non-asymptotic bounds in a streamlined manner. Furthermore, our discretisations differ from the \gls*{mpgd} routine: the \gls*{kiplmc1} algorithm does not use a gradient correction step, whilst \gls*{kiplmc2} is based on a different integrator. In Sec.~\ref{sec:experiments}, we show that these algorithms obtain similar stability, while having theoretical guarantees. 

The paper is structured as follows. In Section~\ref{sec:background}, we present the background work on underdamped Langevin diffusions for sampling and interacting particle systems for the \gls*{mmle} problem. Following this introduction, in Section~\ref{sec:algorithms}, we present a new diffusion targeting the solution of the \gls*{mmle} problem and present the associated algorithms \gls*{kiplmc1} and \gls*{kiplmc2}. Section~\ref{sec:nonasymp} presents the nonasymptotic analysis of \gls*{kiplmc} methods, providing a clear evidence of improved theoretical bounds due to the acceleration. Finally, in Section~\ref{sec:experiments}, we present a variety of experiments to support our theoretical findings.

\subsubsection*{Notation}
Let $\mathcal{P}(\mathbb{R}^d)$, for $d\geq 1$, denote the set of probability measures on $(\mathbb{R}^d,\mathcal{B}(\mathbb{R}^d))$, where $\mathcal{B}(\mathbb{R}^d)$ is the Borel $\sigma$-algebra. We write $\langle\cdot,\cdot\rangle$ and $\|\cdot\|$ for the Euclidean inner product and norm on $\mathbb{R}^d$. We also use $[N]=\{1,\dots,N\}$ and $\bN$ for the positive integers. For $p\geq 1$, define the Wasserstein-$p$ distance by
\[
W_p(\pi,\nu) = \left(\inf_{\Gamma\in \mathbf{T}(\pi, \nu)} \int_{\mathbb{R}^d\times\mathbb{R}^d} \|x-y\|^p\,\mathrm{d}\Gamma(x,y)\right)^{1/p},
\]
where $\mathbf{T}(\pi,\nu)$ denotes the set of couplings of $\pi$ and $\nu$, i.e., probability measures $\Gamma$ on $\mathbb{R}^d\times\mathbb{R}^d$ with marginals $\pi$ and $\nu$. We also define $\widetilde{\mathcal{O}}$ as the big $\mathcal{O}$ notation, up to logarithmic orders.
